\import{function.tex}

\section{Partial maps and their natural domains}

Partial maps use the ordinary function graphs. Their compositions retain the
natural domain: only points whose images lie in the next domain can be used.
The underlying domain and image identities already hold for relations.

\begin{theorem}\label{relation_composition_domain}
    Let $R,S$ be relations.
    Then $\dom{S\circ R}=\preimg{R}{\dom{S}}$.
\end{theorem}
\begin{proof}
    Follows by \auto{setext,dom_iff,preimg_iff,circ_iff}.
\end{proof}

\begin{theorem}\label{relation_composition_range}
    Let $R,S$ be relations.
    Then $\ran{S\circ R}=\img{S}{\ran{R}}$.
\end{theorem}
\begin{proof}
    Follows by \auto{setext,ran_iff,img_iff,circ_iff}.
\end{proof}

\begin{theorem}\label{relation_composition_image}
    Let $R,S$ be relations.
    Then $\img{S\circ R}{A}=\img{S}{\img{R}{A}}$.
\end{theorem}
\begin{proof}
    Follows by \auto{setext,img_iff,circ_iff}.
\end{proof}

\begin{theorem}\label{restriction_intersection}
    Let $R$ be a relation.
    Then $\restrl{\restrl{R}{A}}{B}=\restrl{R}{A\inter B}$.
\end{theorem}
\begin{proof}
    Follows by \auto{setext,restrl,inter}.
\end{proof}

\begin{theorem}\label{restriction_apply_on_domain}
    Let $f$ be a function.
    Let $x\in\dom{f}\inter A$.
    Then $(\restrl{f}{A})(x)=f(x)$.
\end{theorem}
\begin{proof}
    $x\in\dom{f}$ and $x\in A$ by \obv{inter}.
    $(x,f(x))\in f$ by \obv{function_apply_elim}.
    $(x,f(x))\in\restrl{f}{A}$ by \obv{restrl_iff}.
    $\restrl{f}{A}$ is a function by \obv{restrl_of_function_is_function}.
    Follows by \obv{function_apply_intro}.
\end{proof}

\begin{theorem}\label{restriction_preimage}
    Let $R$ be a relation.
    Then $\preimg{\restrl{R}{A}}{B}=\preimg{R}{B}\inter A$.
\end{theorem}
\begin{proof}
    Follows by \auto{setext,preimg_iff,restrl_iff,inter}.
\end{proof}

\begin{theorem}\label{partial_composition_apply}
    Let $f,g$ be functions.
    Suppose $x\in\dom{f}$ and $f(x)\in\dom{g}$.
    Then $(g\circ f)(x)=g(f(x))$.
\end{theorem}
\begin{proof}
    $(x,f(x))\in f$ and $(f(x),g(f(x)))\in g$ by
        \obv{function_apply_elim}.
    Hence $(x,g(f(x)))\in g\circ f$ by \obv{circ_elem_intro}.
    $g\circ f$ is a function by \obv{function_circ}.
    Follows by \obv{function_apply_intro}.
\end{proof}

\begin{theorem}\label{partial_composition_domain_iff}
    Let $f,g$ be functions.
    Then $x\in\dom{g\circ f}$ iff
        $x\in\dom{f}$ and $f(x)\in\dom{g}$.
\end{theorem}
\begin{proof}
    $\dom{g\circ f}=\preimg{f}{\dom{g}}$ by
        \obv{relation_composition_domain}.
    Follows by \auto{preimg_iff,function_apply_iff}.
\end{proof}

\begin{theorem}\label{threefold_partial_composition_domain}
    Let $f,g,h$ be functions.
    Then $\dom{h\circ(g\circ f)}
        =\preimg{f}{\preimg{g}{\dom{h}}}$.
\end{theorem}
\begin{proof}
    $g\circ f$ is a function by \obv{function_circ}.
    $\dom{h\circ(g\circ f)}=\preimg{g\circ f}{\dom{h}}$ by
        \obv{relation_composition_domain}.
    $\preimg{g\circ f}{\dom{h}}
        =\preimg{f}{\preimg{g}{\dom{h}}}$ by \obv{preimg_circ}.
\end{proof}

\begin{theorem}\label{image_domain_trace}
    Let $R$ be a relation.
    Then $\img{R}{\dom{R}\inter A}=\img{R}{A}$.
\end{theorem}
\begin{proof}
    Follows by \auto{setext,img_iff,dom_iff,inter}.
\end{proof}

\begin{theorem}\label{restriction_range}
    Let $R$ be a relation.
    Then $\ran{\restrl{R}{A}}=\img{R}{A}$.
\end{theorem}
\begin{proof}
    Follows by \auto{setext,ran_iff,restrl_iff,img_iff}.
\end{proof}

\begin{theorem}\label{restriction_exact_type_to_image}
    Let $f$ be a function.
    Suppose $U\subseteq\dom{f}$.
    Then $\restrl{f}{U}\in\funs{U}{\img{f}{U}}$.
\end{theorem}
\begin{proof}
    $\restrl{f}{U}$ is a function by \obv{restrl_of_function_is_function}.
    $\dom{\restrl{f}{U}}=\dom{f}\inter U$ by \obv{dom_of_restrl_eq_inter}.
    $\dom{f}\inter U=U$ by \obv{inter_absorb_supseteq_left}.
    $\ran{\restrl{f}{U}}=\img{f}{U}$ by \obv{restriction_range}.
    Follows by \auto{function_exact_type_to_range}.
\end{proof}

\begin{theorem}\label{image_preimage_range_trace}
    Let $f$ be a function.
    Then $\img{f}{\preimg{f}{B}}=\ran{f}\inter B$.
\end{theorem}
\begin{proof}
    Show that every element of $\img{f}{\preimg{f}{B}}$
        is an element of $\ran{f}\inter B$.
    \begin{subproof}
        Fix $y\in\img{f}{\preimg{f}{B}}$.
        Take $x\in\preimg{f}{B}$ such that $(x,y)\in f$ by \obv{img_iff}.
        Take $z\in B$ such that $(x,z)\in f$ by \obv{preimg_iff}.
        $y=z$ by \obv{rightunique}.
        $y\in\ran{f}$ by \obv{ran_iff}.
        Follows by \obv{inter}.
    \end{subproof}
    Show that every element of $\ran{f}\inter B$
        is an element of $\img{f}{\preimg{f}{B}}$.
    \begin{subproof}
        Fix $y\in\ran{f}\inter B$.
        $y\in\ran{f}$ and $y\in B$ by \obv{inter}.
        Take $x$ such that $(x,y)\in f$ by \obv{ran_iff}.
        $x\in\preimg{f}{B}$ by \auto{preimg_iff}.
        Follows by \auto{img_iff}.
    \end{subproof}
    Follows by \obv{setext}.
\end{proof}

\begin{theorem}\label{preimage_image_domain_trace}
    Let $f$ be an injection.
    Then $\preimg{f}{\img{f}{A}}=\dom{f}\inter A$.
\end{theorem}
\begin{proof}
    $\converse{f}$ is a function by \obv{injective_converse_is_function}.
    $\img{\converse{f}}{\preimg{\converse{f}}{A}}
        =\ran{\converse{f}}\inter A$ by \obv{image_preimage_range_trace}.
    $\preimg{\converse{f}}{A}=\img{\converse{\converse{f}}}{A}$ by
        \obv{preim_eq_img_of_converse}.
    $\converse{\converse{f}}=f$ by \obv{converse_converse_eq}.
    $\ran{\converse{f}}=\dom{f}$ by \obv{ran_converse}.
    $\preimg{f}{\img{f}{A}}=\img{\converse{f}}{\img{f}{A}}$ by
        \obv{preim_eq_img_of_converse}.
\end{proof}

\begin{theorem}\label{injective_image_intersection}
    Let $f$ be an injection.
    Then $\img{f}{A\inter B}=\img{f}{A}\inter\img{f}{B}$.
\end{theorem}
\begin{proof}
    Show that every element of $\img{f}{A\inter B}$
        is an element of $\img{f}{A}\inter\img{f}{B}$.
    \begin{subproof}
        Fix $y\in\img{f}{A\inter B}$.
        Take $x\in A\inter B$ such that $(x,y)\in f$ by \obv{img_iff}.
        $x\in A$ and $x\in B$ by \obv{inter}.
        $y\in\img{f}{A}$ and $y\in\img{f}{B}$ by \obv{img_elem_intro}.
        Follows by \obv{inter}.
    \end{subproof}
    Show that every element of $\img{f}{A}\inter\img{f}{B}$
        is an element of $\img{f}{A\inter B}$.
    \begin{subproof}
        Fix $y\in\img{f}{A}\inter\img{f}{B}$.
        Take $a\in A$ such that $(a,y)\in f$ by \obv{inter,img_iff}.
        Take $b\in B$ such that $(b,y)\in f$ by \obv{inter,img_iff}.
        $a=b$ by \obv{injective}.
        Hence $a\in A\inter B$ by \obv{inter}.
        Follows by \obv{img_elem_intro}.
    \end{subproof}
    Follows by \obv{setext}.
\end{proof}

\begin{theorem}\label{restriction_converse}
    Let $f$ be an injection.
    Then $\converse{\restrl{f}{A}}
        =\restrl{\converse{f}}{\img{f}{A}}$.
\end{theorem}
\begin{proof}
    $\converse{\restrl{f}{A}}$ and $\converse{f}$ are relations by
        \obv{converse_is_relation}.
    $\restrl{\converse{f}}{\img{f}{A}}$ is a relation by
        \obv{relation,restrl_subseteq,subseteq}.
    Follows by \auto{relext,converse_iff,restrl_iff,img_iff,injective}.
\end{proof}

\begin{theorem}\label{partial_composition_restrict_target}
    Let $f,g$ be functions.
    Then $(\restrl{g}{B})\circ f
        =\restrl{g\circ f}{\preimg{f}{B}}$.
\end{theorem}
\begin{proof}
    $(\restrl{g}{B})\circ f$ and $g\circ f$ are functions by
        \obv{restrl_of_function_is_function,function_circ}.
    $\restrl{g\circ f}{\preimg{f}{B}}$ is a function by
        \obv{restrl_of_function_is_function}.
    Follows by
        \auto{relext,circ_iff,restrl_iff,preimg_iff,rightunique}.
\end{proof}

\begin{theorem}\label{partial_composition_local_restriction}
    Let $f,g$ be functions.
    Suppose $W\subseteq\preimg{f}{V}$.
    Then $\restrl{g\circ f}{W}=(\restrl{g}{V})\circ(\restrl{f}{W})$.
\end{theorem}
\begin{proof}
    $(\restrl{g}{V})\circ f=\restrl{g\circ f}{\preimg{f}{V}}$ by
        \obv{partial_composition_restrict_target}.
    $(\restrl{g}{V})\circ(\restrl{f}{W})
        =\restrl{(\restrl{g}{V})\circ f}{W}$ by \obv{restrl_circ}.
    $g\circ f$ is a relation by \obv{circ_is_relation}.
    $\restrl{\restrl{g\circ f}{\preimg{f}{V}}}{W}
        =\restrl{g\circ f}{\preimg{f}{V}\inter W}$ by
        \obv{restriction_intersection}.
    $\preimg{f}{V}\inter W=W$ by \obv{inter_absorb_supseteq_left}.
    Follows by \auto.
\end{proof}

\begin{theorem}\label{composition_restrict_intermediate_range}
    Let $R,S$ be relations.
    Then $(\restrl{S}{\ran{R}})\circ R=S\circ R$.
\end{theorem}
\begin{proof}
    Follows by \auto{relext,circ_is_relation,circ_iff,restrl_iff,ran_iff}.
\end{proof}

\begin{theorem}\label{transition_restrict_overlap}
    Let $c,d$ be relations.
    Then $d\circ\converse{c}
        =(\restrl{d}{\dom{c}\inter\dom{d}})
            \circ\converse{\restrl{c}{\dom{c}\inter\dom{d}}}$.
\end{theorem}
\begin{proof}
    Let $A=\dom{c}\inter\dom{d}$.
    Let $R=d\circ\converse{c}$.
    Let $S=(\restrl{d}{A})\circ\converse{\restrl{c}{A}}$.
    $R$ and $S$ are relations by \auto{circ_is_relation}.
    Show that for all $x,y$ we have $(x,y)\in R$ iff $(x,y)\in S$.
    \begin{subproof}
        Fix $x,y$.
        Show that if $(x,y)\in R$, then $(x,y)\in S$.
        \begin{subproof}
            Assume $(x,y)\in R$.
            Take $z$ such that $(z,x)\in c$ and $(z,y)\in d$ by
                \auto{circ_iff,converse_iff}.
            $z\in A$ by \auto{dom_iff,inter}.
            $(z,x)\in\restrl{c}{A}$ and $(z,y)\in\restrl{d}{A}$ by
                \auto{restrl_iff}.
            Follows by \auto{circ_iff,converse_iff}.
        \end{subproof}
        If $(x,y)\in S$, then $(x,y)\in R$ by
            \auto{circ_iff,converse_iff,restrl_iff}.
        Follows by \auto.
    \end{subproof}
    $R=S$ by \obv{relext}.
\end{proof}

\begin{theorem}\label{partial_transition_domain}
    Let $c,d$ be relations.
    Then $\dom{d\circ\converse{c}}
        =\img{c}{\dom{c}\inter\dom{d}}$.
\end{theorem}
\begin{proof}
    $\converse{c}$ is a relation by \obv{converse_is_relation}.
    $\dom{d\circ\converse{c}}=\preimg{\converse{c}}{\dom{d}}$ by
        \obv{relation_composition_domain}.
    $\preimg{\converse{c}}{\dom{d}}
        =\img{\converse{\converse{c}}}{\dom{d}}$ by
        \obv{preim_eq_img_of_converse}.
    $\converse{\converse{c}}=c$ by \obv{converse_converse_eq}.
    $\img{c}{\dom{c}\inter\dom{d}}=\img{c}{\dom{d}}$ by
        \obv{image_domain_trace}.
\end{proof}

\begin{theorem}\label{transition_factorization}
    Let $a,c$ be injections and $d$ be a relation.
    Then $(d\circ\converse{a})\circ(a\circ\converse{c})
        =\restrl{d\circ\converse{c}}{\img{c}{\dom{a}}}$.
\end{theorem}
\begin{proof}
    Let $R=(d\circ\converse{a})\circ(a\circ\converse{c})$.
    Let $S=\restrl{d\circ\converse{c}}{\img{c}{\dom{a}}}$.
    $R$ is a relation by \obv{circ_is_relation}.
    $S$ is a relation by \auto{relation,restrl,circ_is_relation}.
    Show that for all $x,y$ we have $(x,y)\in R$ iff $(x,y)\in S$.
    \begin{subproof}
        Fix $x,y$.
        Show that if $(x,y)\in R$, then $(x,y)\in S$.
        \begin{subproof}
            Assume $(x,y)\in R$.
            Take $z$ such that $(x,z)\in a\circ\converse{c}$ and
                $(z,y)\in d\circ\converse{a}$ by \obv{circ_iff}.
            Take $p$ such that $(p,x)\in c$ and $(p,z)\in a$ by
                \auto{circ_iff,converse_iff}.
            Take $q$ such that $(q,z)\in a$ and $(q,y)\in d$ by
                \auto{circ_iff,converse_iff}.
            $p=q$ by \obv{injective}.
            $p\in\dom{a}$ by \auto{dom_iff}.
            $x\in\img{c}{\dom{a}}$ by \auto{img_iff}.
            $(x,y)\in d\circ\converse{c}$ by \auto{circ_iff,converse_iff}.
            Follows by \obv{restrl_iff}.
        \end{subproof}
        Show that if $(x,y)\in S$, then $(x,y)\in R$.
        \begin{subproof}
            Assume $(x,y)\in S$.
            $(x,y)\in d\circ\converse{c}$ and $x\in\img{c}{\dom{a}}$ by
                \obv{restrl_iff}.
            Take $p$ such that $(p,x)\in c$ and $(p,y)\in d$ by
                \auto{circ_iff,converse_iff}.
            Take $q\in\dom{a}$ such that $(q,x)\in c$ by \obv{img_iff}.
            $p=q$ by \obv{injective}.
            Take $z$ such that $(p,z)\in a$ by \auto{dom_iff}.
            $(x,z)\in a\circ\converse{c}$ and $(z,y)\in d\circ\converse{a}$ by
                \auto{circ_iff,converse_iff}.
            Follows by \auto{circ_iff}.
        \end{subproof}
        Follows by \auto.
    \end{subproof}
    $R=S$ by \obv{relext}.
\end{proof}

\begin{theorem}\label{restricted_threefold_partial_composition_domain}
    Let $f,g,h$ be functions.
    Then $\dom{h\circ(g\circ\restrl{f}{A})}
        =\preimg{f}{\preimg{g}{\dom{h}}}\inter A$.
\end{theorem}
\begin{proof}
    $g\circ f$ and $h\circ(g\circ f)$ are functions by
        \obv{function_circ}.
    $\restrl{g\circ f}{A}=g\circ\restrl{f}{A}$ by
        \obv{restrl_circ}.
    $\restrl{h\circ(g\circ f)}{A}
        =h\circ\restrl{g\circ f}{A}$ by \obv{restrl_circ}.
    $\dom{\restrl{h\circ(g\circ f)}{A}}
        =\dom{h\circ(g\circ f)}\inter A$ by
        \obv{dom_of_restrl_eq_inter}.
    $\dom{h\circ(g\circ f)}
        =\preimg{f}{\preimg{g}{\dom{h}}}$ by
        \obv{threefold_partial_composition_domain}.
\end{proof}

\subsection{Identities and change of source}

\begin{theorem}\label{relation_circ_identity}
    Let $R$ be a relation.
    Suppose $\dom{R}\subseteq A$.
    Then $R\circ\identity{A}=R$.
\end{theorem}
\begin{proof}
    For all $x,y$ we have $(x,y)\in R\circ\identity{A}$ iff $(x,y)\in R$ by
        \auto{circ_iff,id_iff,dom_iff,subseteq}.
    Follows by \obv{relext,circ_is_relation}.
\end{proof}

\begin{theorem}\label{relation_identity_circ}
    Let $R$ be a relation.
    Suppose $\ran{R}\subseteq A$.
    Then $\identity{A}\circ R=R$.
\end{theorem}
\begin{proof}
    For all $x,y$ we have $(x,y)\in\identity{A}\circ R$ iff $(x,y)\in R$ by
        \auto{circ_iff,id_iff,ran_iff,subseteq}.
    Follows by \obv{relext,circ_is_relation}.
\end{proof}

\begin{theorem}\label{composition_bijection_roundtrip}
    Let $R$ be a relation and $h\in\bijections{A}{B}$.
    Suppose $\dom{R}\subseteq A$.
    Then $(R\circ\converse{h})\circ h=R$.
\end{theorem}
\begin{proof}
    $\converse{h}\circ h=\identity{A}$ by \obv{bijection_converse_circ}.
    $R\circ\identity{A}=R$ by \obv{relation_circ_identity}.
    $(R\circ\converse{h})\circ h=R\circ(\converse{h}\circ h)$ by \obv{circ_assoc}.
\end{proof}

\begin{theorem}\label{identity_restriction}
    Then $\restrl{\identity{A}}{B}=\identity{A\inter B}$.
\end{theorem}
\begin{proof}
    $\identity{A}$ is a relation by \obv{id_is_relation}.
    For all $x,y$ we have $(x,y)\in\restrl{\identity{A}}{B}$
        iff $(x,y)\in\identity{A\inter B}$ by \auto{restrl_iff,id_iff,inter}.
    Follows by \obv{relext,id_is_relation}.
\end{proof}

\begin{theorem}\label{identity_partial_composition}
    Then $\identity{B}\circ\identity{A}=\identity{A\inter B}$.
\end{theorem}
\begin{proof}
    Follows by
        \auto{relext,circ_is_relation,id_is_relation,circ_iff,id_iff,inter}.
\end{proof}

\begin{corollary}\label{disjoint_identity_composition_domain}
    Suppose $A\inter B=\emptyset$.
    Then $\dom{\identity{B}\circ\identity{A}}=\emptyset$.
\end{corollary}
\begin{proof}
    $\identity{B}\circ\identity{A}=\identity{A\inter B}$ by
        \obv{identity_partial_composition}.
    $\dom{\identity{A\inter B}}=A\inter B$ by \obv{id_dom}.
\end{proof}
